% !TeX root = ../教材题库_全部条目_转写工作稿.tex
% Source: XJ-B1-C06-HW-P256-Y0260
% 第6章 / 6.4 用样本估计总体 / 教材 p.256 / PDF 第263页
% OCR未校对：公式与文字均可能有误；未完成数学TeX转写；禁止直接发给学生
% 原题及相邻原解截图（核对用）：../题目原图/XJ-B1-C06-HW-P256-Y0260-01.webp
% 本文件只在转写完成后作为题干使用；OCR 条目尚不是可靠题干。
\begin{bankitem}{习题6.4 第9题}{第6章 / 6.4 用样本估计总体 / 教材 p.256 / PDF 第263页}{XJ-B1-C06-HW-P256-Y0260}{OCR未校对：公式与文字均可能有误；未完成数学TeX转写；禁止直接发给学生}
\begin{ocrdraft}
[教材p.256 / PDF第263页]
\par
9．人体测量的数据以第百分位数（记为撼）作为一种指标界值．最常用的是
\par
P5，P50,P!〕5三种．在身高中，我们称P5为矮身材，P0为中身材，P9为高身材．
\par
现调查得到如下所示的20名19岁中国女性的身高数据（单位：
\par
巧2
\par
巧2153154155156
\par
158
\par
159
\par
160
\par
161
\par
162
\par
162163163165167
\par
168
\par
170
\par
171
\par
172
\par
试分别求矮身材、中身材、高身材的界值．
\end{ocrdraft}
\end{bankitem}
